\scanpage{21}
% sheet: reused_typed
% hand: typed
% archive sheet number: 20
% @ text / head, typescript pagination
\mspage{- 65 -}
% @ text / upper third, Proposition 2
PROPOSITION 2 (lemme de Schwarz).-- \uline{Soit} $f$ \uline{une fonction holomorphe dans un domaine} $A\handwritten{\subset\unsure{\mathbb C|C}}$ \uline{contenant le disque} $|z|\handwritten{\leq}1$ , \uline{et supposons que} $f$ \uline{ait un zéro d'ordre $k$ au point} $z=0$ . \uline{Dans ces conditions}, \uline{si} $M$ \uline{est le maximum de} $\|f(z)\|$ \uline{pour} $|z|=1$ , \uline{on a} $\|f(z)\|\handwritten{\leq}M|z|^k$ \uline{pour tout} $z$ \uline{tel que} $|z|\handwritten{\leq}1$ .

% @ text / upper third, proof
Il suffit de remarquer que $z^{-k}f(z)$ est holomorphe dans $A$ et que le maximum de $\|z\unsure{\vert}^{-k}f(z)\|$ pour $|z|=1$ est égal à $M$ .

% @ text / middle, section heading
\uline{6. Ensembles équicontinus de fonctions holomorphes}.

% @ text / middle, Proposition 3
PROPOSITION 3.-- \uline{Soit} $A\handwritten{\subset\mathbb C}^p$ \uline{un domaine}, $\handwritten{\Phi}$ \uline{un ensemble de fonctions holomorphes dans} $A$ , \uline{à valeurs dans un espace de Banach} $E$ . \uline{Si les fonctions de} $\handwritten{\Phi}$ \uline{sont uniformément bornées dans toute partie compacte de} $A$ , $\handwritten{\Phi}$ \uline{est un ensemble équicontinu dans} $A$ .

% @ text / middle to lower third, proof
En effet, soit $\handwritten{a}$ un point de $A$ , $P$ un polycylindre fermé de centre $\handwritten{a}$ et de rayons $\handwritten{\rho}$ , contenu dans $P$ . Par hypothèse, il existe un nombre $M$ tel que $\|f(\handwritten{z})\|\handwritten{\leq}M$ pour tout $\handwritten{z\in}P$ . Soit $Q$ le polycylindre fermé de centre $\handwritten{a}$ et de rayons $\handwritten{\rho}/2$ ; pour tout $\handwritten{x\in}Q$ , le polycylindre fermé de centre $\handwritten{x}$ et de rayons $\handwritten{\rho}/2$ est contenu dans $P$ , donc la formule (5) du $\S\,3$ donne en particulier les majorations

% @ display / lower middle, (3)
\[
\lVert\handwritten{\partial}f/\handwritten{\partial}x_k\rVert\handwritten{\leq}2M/\handwritten{\rho}\qquad ,\tag*{\makebox[0pt][r]{\makebox[\displaywidth][l]{(3)}}}
\]
% @ text / after (3)
valables pour tout $\handwritten{x\in}Q$ et toute fonction $f\handwritten{\in\Phi}$ . Or, on peut écrire, pour tout $\handwritten{y\in}Q$

% @ display / lower third, integral
\[
f(\handwritten{y})-f(\handwritten{a})=\handwritten{\int_{\sigma}}\bigl(\handwritten{\sum_k}\ \frac{\handwritten{\partial} f}{\handwritten{\partial} x_k}\,dx_k\bigr)
\]
% @ text / lower third, parametrization
l'intégrale étant prise le long du segment $\handwritten{\sigma}:t\handwritten{\longrightarrow a}+t(\handwritten{y}-\handwritten{a})$, avec $0\handwritten{\leq}t\handwritten{\leq}1$ ; on en déduit immédiatement, compte-tenu de (3), que

% @ display / last inequality
\[
\lVert f(\handwritten{y})-f(\handwritten{a})\rVert\handwritten{\leq}\frac{2pM}{\handwritten{\rho}}\lVert\handwritten{y}-\handwritten{a}\rVert
\]
% @ text / foot
pour \uline{toute} fonction $f\handwritten{\in\Phi}$ et tout $\handwritten{y\in}Q$ , d'où la proposition.
